\chapter{Monitors and detectors}
\index{Library!Components!monitors}
\index{Monitors|textbf}

In real neutron experiments, detectors and monitors play quite
different roles. One wants the detectors to be as efficient as
possible, counting all neutrons (absorbing them in the process),
while the monitors measure the intensity of the incoming beam, and must
as such be almost transparent, interacting only with (roughly) 0.1-1\%
of the neutrons passing by. In computer simulations, it is
of course possible to detect every neutron ray without
absorbing it or disturbing any of its parameters. Hence, the two components
have very similar functions in the simulations, and we do
not distinguish between them. For simplicity, they are from here on
just called \textbf{monitors}.

Another important difference between computer simulations
and real experiments is
that one may allow the monitor to be sensitive to any neutron property,
as {\em e.g.} direction, energy, and divergence, in addition to what
is found in real-world detectors (space and time). One may, in
fact, let the monitor record correlations between these properties,
as seen for example in the divergence/position sensitive monitor in
section~\ref{s:divpos-monitor}.

When a monitor detects a neutron ray,
a number counting variable is incremented: $n_i = n_{i-1}+1$.
In addition, the neutron
weight $p_i$ is added to the weight counting variable:
$I_i = I_{i-1} + p_i$,
and the second moment of the weight is
updated: $M_{2,i} = M_{2,i-1} + p_i^2$.
As also discussed chapter \ref{s:MCtechniques}, after a simulation of $N$ rays
the detected intensity (in units of neutrons/sec.) is $I_N$,
while the estimated errorbar is $\sqrt{M_{2,N}^2}$.

Many different monitor components have been developed for
\MCS , but we have decided to support only the most important ones.
One example of the monitors we have omitted is the single monitor,
\textbf{Monitor},
that measures just one number (with errorbars) per simulation.
This effect is mirrored by any of the 1- or 2-dimensional components
we support, e.g. the \textrm{PSD\_monitor}.
In case additional functionality of monitors is required,
the few code lines of existing monitors can easily be modified.

However, the ultimate solution is the use of the
``Swiss army knife'' of monitors, \textbf{Monitor\_nD}, that can face
almost any simulation requirement,
but will prove challenging for users who like to perform own modifications.
When compiling an instrument for GPU/OpenACC (see the User Manual, section
on GPU acceleration), \texttt{Monitor\_nD} automatically falls back to the
companion component \textbf{Monitor\_nD\_noacc} for any \texttt{options}
string not yet supported on GPU, via the \texttt{NOACC}/\texttt{CPU}
funneling mechanism (see the User Manual's kernel chapter); this is
handled transparently and does not normally require any instrument change.

The monitors presented in detail below are a representative selection;
the library additionally contains cylindrical and spherical variants of
most of them (e.g. \texttt{Cyl\_monitor}, \texttt{PSD\_monitor\_4PI}),
energy-transfer and $S(q,\omega)$-aware monitors
(\texttt{Sqw\_monitor}, \texttt{TOFLambda\_monitor}), polarisation-aware
monitors (\texttt{Pol\_monitor}, \texttt{PolLambda\_monitor}), and
flexible, user-configurable histogrammers (\texttt{Flex\_monitor\_1D/2D/3D}).
Use \verb+mcdoc+ (User Manual, section on McDoc) for the complete,
current list.

\newpage
\section{The \texttt{TOF\_monitor} McStas Component}
Rectangular Time-of-flight monitor.

\subsection*{Identification}
\begin{itemize}
  \item \textbf{Site:} 
  \item \textbf{Author:} KN, M. Hagen
  \item \textbf{Origin:} Risoe
  \item \textbf{Date:} August 1998
\end{itemize}

\subsection*{Description}
\begin{lstlisting}

\end{lstlisting}

\subsection*{Input parameters}
Parameters in \textbf{boldface} are required; the others are optional.

\begin{longtable}{p{0.22\textwidth}p{0.12\textwidth}p{0.46\textwidth}p{0.14\textwidth}}
\toprule
\textbf{Name} & \textbf{Unit} & \textbf{Description} & \textbf{Default} \\
\midrule
\endhead
nt & 1 & Number of time bins & 20 \\
filename & string & Name of file in which to store the detector image & 0 \\
xmin & m & Lower x bound of detector opening & -0.05 \\
xmax & m & Upper x bound of detector opening & 0.05 \\
ymin & m & Lower y bound of detector opening & -0.05 \\
ymax & m & Upper y bound of detector opening & 0.05 \\
xwidth & m & Width of detector. Overrides xmin, xmax & 0 \\
yheight & m & Height of detector. Overrides ymin, ymax & 0 \\
tmin & mu-s & Lower time limit & 0 \\
tmax & mu-s & Upper time limit & 0 \\
dt & mu-s & Length of each time bin & 1.0 \\
restore\_neutron & 1 & If set, the monitor does not influence the neutron state & 0 \\
nowritefile & 1 & If set, monitor will skip writing to disk & 0 \\
\bottomrule
\end{longtable}

\subsection*{Links}
\begin{itemize}
  \item \href{run:/home/willend/willend-McCode/mcstas-comps/monitors/TOF_monitor.comp}{Source code} for \texttt{TOF\_monitor.comp}.
\end{itemize}
\IfFileExists{TOF_monitor_static.tex}{\input{TOF_monitor_static.tex}}{}

\section{The \texttt{TOF2E\_monitor} McStas Component}
TOF-sensitive monitor, converting to energy

\subsection*{Identification}
\begin{itemize}
  \item \textbf{Site:} 
  \item \textbf{Author:} Kim Lefmann and Helmuth Schoeber
  \item \textbf{Origin:} Risoe
  \item \textbf{Date:} Sept. 13, 2006
\end{itemize}

\subsection*{Description}
\begin{lstlisting}
A square single monitor that measures the energy of the incoming neutrons
from their time-of-flight

Example: TOF2E_monitor(xmin=-0.1, xmax=0.1, ymin=-0.1, ymax=0.1, Emin=1, Emax=50, nE=20, filename="Output.nrj", L_flight=20, T_zero=0)
\end{lstlisting}

\subsection*{Input parameters}
Parameters in \textbf{boldface} are required; the others are optional.

\begin{longtable}{p{0.22\textwidth}p{0.12\textwidth}p{0.46\textwidth}p{0.14\textwidth}}
\toprule
\textbf{Name} & \textbf{Unit} & \textbf{Description} & \textbf{Default} \\
\midrule
\endhead
nE & 1 & Number of energy channels & 20 \\
filename & string & Name of file in which to store the detector image & 0 \\
nowritefile & 1 & If set, monitor will skip writing to disk & 0 \\
xmin & m & Lower x bound of detector opening & -0.05 \\
xmax & m & Upper x bound of detector opening & 0.05 \\
ymin & m & Lower y bound of detector opening & -0.05 \\
ymax & m & Upper y bound of detector opening & 0.05 \\
xwidth & m & Width of detector. Overrides xmin, xmax & 0 \\
yheight & m & Height of detector. Overrides ymin, ymax & 0 \\
\textbf{Emin} & meV & Minimum energy to detect &  \\
\textbf{Emax} & meV & Maximum energy to detect &  \\
\textbf{T\_zero} & s & Zero point in time &  \\
\textbf{L\_flight} & m & flight length user in conversion &  \\
restore\_neutron & 1 & If set, the monitor does not influence the neutron state & 0 \\
\bottomrule
\end{longtable}

\subsection*{Links}
\begin{itemize}
  \item \href{run:/home/willend/willend-McCode/mcstas-comps/monitors/TOF2E_monitor.comp}{Source code} for \texttt{TOF2E\_monitor.comp}.
\end{itemize}
\IfFileExists{TOF2E_monitor_static.tex}{\input{TOF2E_monitor_static.tex}}{}

\section{The \texttt{E\_monitor} McStas Component}
Energy-sensitive monitor.

\subsection*{Identification}
\begin{itemize}
  \item \textbf{Site:} 
  \item \textbf{Author:} Kristian Nielsen and Kim Lefmann
  \item \textbf{Origin:} Risoe
  \item \textbf{Date:} April 20, 1998
\end{itemize}

\subsection*{Description}
\begin{lstlisting}
A square single monitor that measures the energy of the incoming neutrons.

Example: E_monitor(xmin=-0.1, xmax=0.1, ymin=-0.1, ymax=0.1,
Emin=1, Emax=50, nE=20, filename="Output.nrj")
\end{lstlisting}

\subsection*{Input parameters}
Parameters in \textbf{boldface} are required; the others are optional.

\begin{longtable}{p{0.22\textwidth}p{0.12\textwidth}p{0.46\textwidth}p{0.14\textwidth}}
\toprule
\textbf{Name} & \textbf{Unit} & \textbf{Description} & \textbf{Default} \\
\midrule
\endhead
nE & 1 & Number of energy channels & 20 \\
filename & string & Name of file in which to store the detector image & 0 \\
xmin & m & Lower x bound of detector opening & -0.05 \\
xmax & m & Upper x bound of detector opening & 0.05 \\
ymin & m & Lower y bound of detector opening & -0.05 \\
ymax & m & Upper y bound of detector opening & 0.05 \\
nowritefile & 1 & If set, monitor will skip writing to disk & 0 \\
xwidth & m & Width of detector. Overrides xmin, xmax. & 0 \\
yheight & m & Height of detector. Overrides ymin, ymax. & 0 \\
\textbf{Emin} & meV & Minimum energy to detect &  \\
\textbf{Emax} & meV & Maximum energy to detect &  \\
restore\_neutron & 1 & If set, the monitor does not influence the neutron state & 0 \\
\bottomrule
\end{longtable}

\subsection*{Links}
\begin{itemize}
  \item \href{run:/home/willend/willend-McCode/mcstas-comps/monitors/E_monitor.comp}{Source code} for \texttt{E\_monitor.comp}.
\end{itemize}
\IfFileExists{E_monitor_static.tex}{\input{E_monitor_static.tex}}{}

\section{The \texttt{L\_monitor} McStas Component}
Wavelength-sensitive monitor.

\subsection*{Identification}
\begin{itemize}
  \item \textbf{Site:} 
  \item \textbf{Author:} Kristian Nielsen and Kim Lefmann
  \item \textbf{Origin:} Risoe
  \item \textbf{Date:} April 20, 1998
\end{itemize}

\subsection*{Description}
\begin{lstlisting}
A square single monitor that measures the wavelength of the incoming
neutrons.

Example: L_monitor(xmin=-0.1, xmax=0.1, ymin=-0.1, ymax=0.1, nL=20, filename="Output.L", Lmin=2, Lmax=10)
\end{lstlisting}

\subsection*{Input parameters}
Parameters in \textbf{boldface} are required; the others are optional.

\begin{longtable}{p{0.22\textwidth}p{0.12\textwidth}p{0.46\textwidth}p{0.14\textwidth}}
\toprule
\textbf{Name} & \textbf{Unit} & \textbf{Description} & \textbf{Default} \\
\midrule
\endhead
nL & 1 & Number of wavelength channels & 20 \\
filename & string & Name of file in which to store the detector image & 0 \\
nowritefile & 1 & If set, monitor will skip writing to disk & 0 \\
xmin & m & Lower x bound of detector opening & -0.05 \\
xmax & m & Upper x bound of detector opening & 0.05 \\
ymin & m & Lower y bound of detector opening & -0.05 \\
ymax & m & Upper y bound of detector opening & 0.05 \\
xwidth & m & Width of detector. Overrides xmin, xmax. & 0 \\
yheight & m & Height of detector. Overrides ymin, ymax. & 0 \\
\textbf{Lmin} & AA & Minimum wavelength to detect &  \\
\textbf{Lmax} & AA & Maximum wavelength to detect &  \\
restore\_neutron & 1 & If set, the monitor does not influence the neutron state & 0 \\
\bottomrule
\end{longtable}

\subsection*{Links}
\begin{itemize}
  \item \href{run:/home/willend/willend-McCode/mcstas-comps/monitors/L_monitor.comp}{Source code} for \texttt{L\_monitor.comp}.
\end{itemize}
\IfFileExists{L_monitor_static.tex}{\input{L_monitor_static.tex}}{}

\input{monitors/PSD_monitor}

\section{The \texttt{Divergence\_monitor} McStas Component}
Horizontal+vertical divergence monitor.

\subsection*{Identification}
\begin{itemize}
  \item \textbf{Site:} 
  \item \textbf{Author:} Kim Lefmann
  \item \textbf{Origin:} Risoe
  \item \textbf{Date:} Nov. 11, 1998
\end{itemize}

\subsection*{Description}
\begin{lstlisting}
A divergence sensitive monitor. The counts are distributed in
(n times m) pixels.

Example: Divergence_monitor(nh=20, nv=20, filename="Output.pos",
xmin=-0.1, xmax=0.1, ymin=-0.1, ymax=0.1,
maxdiv_h=2, maxdiv_v=2)
\end{lstlisting}

\subsection*{Input parameters}
Parameters in \textbf{boldface} are required; the others are optional.

\begin{longtable}{p{0.22\textwidth}p{0.12\textwidth}p{0.46\textwidth}p{0.14\textwidth}}
\toprule
\textbf{Name} & \textbf{Unit} & \textbf{Description} & \textbf{Default} \\
\midrule
\endhead
nh & 1 & Number of pixel rows & 20 \\
nv & 1 & Number of pixel columns & 20 \\
filename & str & Name of file in which to store the detector image text & 0 \\
xmin & m & Lower x bound of detector opening & -0.05 \\
xmax & m & Upper x bound of detector opening & 0.05 \\
ymin & m & Lower y bound of detector opening & -0.05 \\
ymax & m & Upper y bound of detector opening & 0.05 \\
nowritefile & 1 & If set, monitor will skip writing to disk & 0 \\
xwidth & m & Width of detector. Overrides xmin, xmax & 0 \\
yheight & m & Height of detector. Overrides ymin, ymax & 0 \\
maxdiv\_h & degrees & Maximal vertical divergence detected & 2 \\
maxdiv\_v & degrees & Maximal vertical divergence detected & 2 \\
restore\_neutron & 1 & If set, the monitor does not influence the neutron state & 0 \\
nx & 1 &  & 0 \\
ny & 1 & Vector definition of "forward" direction wrt. divergence, to be used e.g. when the monitor is rotated into the horizontal plane & 0 \\
nz & 1 &  & 1 \\
\bottomrule
\end{longtable}

\subsection*{Links}
\begin{itemize}
  \item \href{run:/home/willend/willend-McCode/mcstas-comps/monitors/Divergence_monitor.comp}{Source code} for \texttt{Divergence\_monitor.comp}.
\end{itemize}
\IfFileExists{Divergence_monitor_static.tex}{\input{Divergence_monitor_static.tex}}{}

\section{The \texttt{DivPos\_monitor} McStas Component}
Divergence/position monitor (acceptance diagram).

\subsection*{Identification}
\begin{itemize}
  \item \textbf{Site:} 
  \item \textbf{Author:} Kristian Nielsen
  \item \textbf{Origin:} Risoe
  \item \textbf{Date:} 1999
\end{itemize}

\subsection*{Description}
\begin{lstlisting}
2D detector for intensity as a function of position
and divergence, either horizontally or vertically (depending on the flag vertical).
This gives information similar to an aceptance diagram used
eg. to investigate beam profiles in neutron guides.

Example: DivPos_monitor(nh=20, ndiv=20, filename="Output.dip",
xmin=-0.1, xmax=0.1, ymin=-0.1, ymax=0.1, maxdiv_h=2)
\end{lstlisting}

\subsection*{Input parameters}
Parameters in \textbf{boldface} are required; the others are optional.

\begin{longtable}{p{0.22\textwidth}p{0.12\textwidth}p{0.46\textwidth}p{0.14\textwidth}}
\toprule
\textbf{Name} & \textbf{Unit} & \textbf{Description} & \textbf{Default} \\
\midrule
\endhead
nb & 1 & Number of bins in position. & 20 \\
ndiv & 1 & Number of bins in divergence. & 20 \\
filename & string & Name of file in which to store the detector image. & 0 \\
xmin & m & Lower x bound of detector opening & -0.05 \\
xmax & m & Upper x bound of detector opening & 0.05 \\
ymin & m & Lower y bound of detector opening & -0.05 \\
ymax & m & Upper y bound of detector opening & 0.05 \\
xwidth & m & Width of detector. Overrides xmin,xmax. & 0 \\
yheight & m & Height of detector. Overrides ymin,ymax. & 0 \\
maxdiv & degrees & Maximal divergence detected. & 2 \\
restore\_neutron & 1 & If set, the monitor does not influence the neutron state. & 0 \\
nx & 1 &  & 0 \\
ny & 1 & Vector definition of "forward" direction wrt. divergence, to be used e.g. when the monitor is rotated into the horizontal plane & 0 \\
nz & 1 &  & 1 \\
vertical & 1 & Monitor intensity as a function of vertical divergence and position. & 0 \\
nowritefile & 1 & If set, monitor will skip writing to disk & 0 \\
\bottomrule
\end{longtable}

\subsection*{Links}
\begin{itemize}
  \item \href{run:/home/willend/willend-McCode/mcstas-comps/monitors/DivPos_monitor.comp}{Source code} for \texttt{DivPos\_monitor.comp}.
\end{itemize}
\IfFileExists{DivPos_monitor_static.tex}{\input{DivPos_monitor_static.tex}}{}

\newpage
\input{monitors/Monitor_nD}

\newpage
\input{monitors/Res_monitor}
